\sgd{Save three paragraphs (hopefully).}
We are witnessing a technological paradigm shift to move away from service centralization, a hallmark of the cloud computing and web services industries in the past two decades.
This \emph{decentralization} movement is catalyzed by the growing concern over content censorship~\cite{twitter_censorship}, data misuse~\cite{fb_misuse}, monopolistic practices~\cite{google_antitrust}, and single point of organization failure~\cite{ftx_debacle}.
Following the wild success of the Bitcoin~\cite{bitcoin} and Ethereum~\cite{ethereum} blockchains, a burgeoning of decentralized services are being deployed, including exchanges~\cite{uniswap}, content delivery network~\cite{meson}, domain name service~\cite{handshake}, messaging~\cite{matrix}, and computation~\cite{dfinity}.

An attractive avenue for decentralization is data storage~\cite{ipfs, filecoin, ar, storj}.
Compared to centrally managed systems~\cite{amazon_s3, google_cloudstore, dynamo}, decentralized storage is better positioned to address the issues of data breaches, censorship, content availability, and high storage cost.
In particular, by storing data across administrative domains, decentralization theoretically can provide storage lifetime longer than any single organization.
We have already seen successful large deployment of decentralized storage systems:
The InterPlanetary File System (IPFS)~\cite{ipfs} has over 300 thousand active participating peers globally, collectively storing more than 1000 PB of data.

A critical requirement for any storage system is \emph{data durability}, \ie, any data stored in the system can be eventually recovered.
To handle frequent server and device failures in large distributed deployments, replication~\cite{smr,pb,pbft} and erasure coding~\cite{reed-solomon} techniques have become an integral part of these systems.
Both techniques use \emph{redundancy} to tolerate failures of individual components, an effective approach when enough copies or encoding symbols survive.
\sgd{Remove one following sentence.}
They work well in centrally managed systems, where central operators have full knowledge and control over the failure domains and placement groups to minimize correlated failures.

\begin{outline}
The most critical design difficulty of the redundancy approach is that the redundancy is only effective if it is distributed across failure domains.
In another word, if there's considerable probability for storage nodes who store the fragments for the same data object to correlated fail in a short interval, data loss may happen even in the presence of redundancy.
\sgd{Just sufficient portion of those storage nodes fail can cause data loss but the sentence is already too complex.} 

The centrally managed systems have well defined/determined \emph{objective} failure domains, \eg geological locations, hardware/software diversities, etc.
But on the other hand, the whole systems are under the same \emph{administrative} failure domain, \ie it's up to system administrators' wills to whether ensure durability of a particular data object.
Data objects can never be distributed beyond single administrative failure domain.

Recent works of decentralized storage \eg Filecoin improves on this with an open, permissionless network that allows any organization to take part in storing the data, so that any single organization is not enough to cause data loss.
However, the solutions are limited in the sense that any data object is only stored by only few organizations.
\sgd{And the organizations are fixed during the whole data object lifetime.}
The situation is not fundamentally changed, and the solutions cannot benefit from more organizations \ie the durability is not further improved when more organizations join the systems.

We argue that in a permissionless system there's no practical way to reliably collect trusted information of administrative failure domains a priori \ie before those domains fail.
\sgd{Mentioning targeted attack here?}
So we cannot deal with administrative failure domains in the same way as dealing with objective ones, but can only ``do the BFT protocol way''.
\sgd{Should carefully describe what is that way.}

We propose a new dimension of measuring the durability of decentralized storage systems: the ability of tolerating faulty organizations.
A durable storage system should ensure data integrity when the faulty organizations control less than certain amount (\eg $1/3$) of total storage capacity.
Filecoin fails to do so.
\sgd{This definition is either nonsense or trivial for centrally managed storage systems.
Those systems have only two cases: faulty organization controls 0 capacity or 100\%.
In the first case they ensure durability while in the second one they do not.}

\sgd{The four above paragraphs somewhat duplicates the centralize/decentralize comparison in the original flow above this outline.
I cannot make them a nature extension on the original flow.
And the original flow cannot be used to attack prior decentralized solutions.
The main difference of these two flows is the original flow makes division on centralized/decentralized systems \ie how many organizations take part in the whole systems, while this outline makes division on how many organizations take part in storing every individual data object.
I feel like do the latter one is always necessary to distinguish \sys from prior decentralized works, and as long as we already do the latter one there's no need to do the former one.}

We propose \sys as a baseline design that can tolerate administrative correlations.
\sys does so by propagating the fragments over the whole system.
The capacity-based virtual node mechanism ensures we can prove the durability.
\sys deploys rateless erasure code along with a set of communication protocols to ensure efficient put and repair operations.

Evaluations.

Contributions.

\sgd{Remove following.}
\end{outline}

A decentralized deployment presents fundamental challenges to the durability guarantees of a storage system.
There is no central authority to manage the behavior and trustworthiness of the participants;
nodes join and leave the system at high churn rate;
there is no reliable source of information regarding configurations and location of participants;
adversarial behaviors are prevalent.
In such open and asynchronous environment, traditional replication, erasure coding, and repair schemes lose their effectiveness in preventing data losses.
Recently, more than 1600 TiB of data on the Filecoin~\cite{filecoin} network lost all replica copies simultaneously~\cite{filecoin_loss}.

In this work, we present a vision to offer \emph{true durability} in decentralized storage.
To do so, we first develop a new fault model that accurately models failure patterns in a decentralized deployment.
The key contribution of our model is \emph{probabilistic failure domains}.
Unlike centrally managed deployments, failure correlation among system participants is not known \emph{a priori}, and is impossible to predict or learn accurately during runtime~\cite{irisstore}.
To account for this uncertainty, our model considers every \emph{combination} of storage devices a failure domain.
The model incorporates a \emph{failure generator} which is a function that assigns a failure rate to each domain.
Crucially, the failure generator, \ie, the exact failure rate of each domain, is unknown to any participants in the system.

Using our newly defined fault model, we develop a new simulation framework to evaluate the durability guarantee of decentralized storage systems.
The simulator takes a failure generator as input.
At every time unit, the simulator samples all failure domains based on the failure rate assigned by the generator.
The union of all storage devices in the sampled failure domains are removed from the simulation, and their stored data are considered lost.
We apply our simulator to evaluate a decentralized storage system that represents fault tolerance approach of prior solution~\cite{irisstore,glacier,cfs,carbonite,ipfs,xfs,oceanstore}.
The system offer strong durability under independent failures.
However, as soon as the generator introduces correlated failures, the system experiences significant data losses.
Increasing the redundancy has marginally benefit beyond a certain point and impacts the overall system storage capacity.

Based on the insights we learned from simulations, we propose a preliminary new decentralized storage design that provides strong durability guarantees.
By expanding the object placement group size to the entire system, our base design can tolerate arbitrary correlated failures within the redundancy limit.
Critically, the larger placement group is not at the expense of reduced storage capacity.
The design additionally applies a \emph{rateless erasure code} to address the fragment uniqueness challenge in a decentralized network.
To improve the practicality of the design, we refine our protocol to use independent sampling to achieve \emph{probabilistic selection} of the placement group.
We also propose \emph{confidential sampling} to defend against active adversaries.

\iffalse
\sgd{The following paragraphs are also new. I think they are formal enough to be called ``content'' instead of ``outline''.}
\lijl{Could add one or two sentences explaining how to (or equivalently, the challenge of) ensure durability in a storage system.}
With our new focusing on the durability, the most critical design consideration becomes how the permissionless storage system react to failures that permanently takes away the fragments stored on the failed storage nodes.
We emphasize on the nature of permissionless systems that, unlike the centralized ones, the correlation among node failures cannot be feasibly reasoned about.
While a centralized storage system can be designed with predetermined failure domains in mind, a permissionless storage system has to be designed against all kinds of correlated failures, and any assumption made to the failure domains may lead to unexpected data loss upon certain correlated failure.
\lijl{We should first educate the readers about failures in a decentralized system.}

\sgd{New.}
We point out that the fault tolerance capability of a permissionless storage system cannot be discussed against a presumed set of failure domains, as a centralized system with planned failure domains.
\lijl{This assertion would require some explanations.
Probably okay after we expand on the failure discussion above.}
Instead, we define the fault tolerance of a permissionless storage system in the notion of \emph{probabilistic failure domains}.
We consider all combinations of storage nodes to be a possible failure domain, and every possible failure domain has its own \emph{failure rate}, \ie, the probability of failing in a time unit.
(For those storage node combinations that never fail together, \ie, do not actually form a failure domain, their failure rates are 0.)
The failure rates of the failure domains are assigned with a function $P(f) = g_f({n_1, n_2, \dots})$ where $n_1, n_2, \dots$ are the nodes of the failure domain.
This function $g_f$ is called \emph{failure generator}.
\lijl{The definition itself is probably not new.
Just that in a centralized system, the failure rate of each possible failure domain is known/controlled by the admin.
So the definition was applied in an implicit way.
We could say that our contribution is formalizing this general theoretical framework for reasoning failure domains (which could be applied to centralized storage), and to define the right failure model under this framework for decentralized storage.
}

\sgd{With this definition we cannot model the failure domain whose failure likelihood changes over time. Be sure of that.}

\sgd{New.}
A permissionless system is in turn modeled as a series of \emph{failure events} that
\begin{itemize}
    \item In every failure event, all storage nodes from a specific failure domain simultaneously fail in the time unit (which is approximated as an instant).
    \item In each time unit, all failure domains \emph{independently} fail according to their failing rates.
    \item Since failure domains overlap with each other, a storage node fails at a time unit if at least one of the failure domains it belongs fail.
    The failed storage nodes permanently leave the system, so each failure domain can only fail once.
    \item The system's \emph{churn rate} is measured as the expected number of failed nodes in every time unit.
    For simplicity, we assume the system remains constant scale, so in every time unit, the number of fresh storage nodes joining the system equals to the number of failed nodes of that time unit.
    \item As the fresh nodes joining, they immediately form new failure domains in combination of all the other nodes.
    The failure rates of the new failure domains are also assigned by the failure generator.
\end{itemize}
With this system model, the faulty degree of the system can be precisely described with the failure generator.
\lijl{This is probably too much detail for intro.
Can summarize the main point in a couple of sentences and elaborate in a later section.}

\sgd{A complete black box failure generator is rigorous but not intuitive. White box it a little if possible.}

\sgd{One more paragraph here to argue we can assume to know failure generator ahead of time, just the actual failure events not predictable.}

\sgd{New.}
With a properly defined fault model we then perform analysis on the prior storage systems.
\lijl{Might need to summarize their durability techniques first.}
In the analysis we simulate a permissionless network with various failure generators.
\lijl{Either here or earlier, we should provide evidence that these failure generators are realistic in a decentralized environment.}
We then populate data objects into the network and observe their persistency over time.
Unsurprisingly, all previous works yield poor mean time to data lost (MTTDL), suggesting they are not able to truthfully tolerate faults.
Further investigation reveals that they are \emph{failure domain sensitive}, meaning their performance \lijl{durability guarantees?} may dramatically degrade because of the presence of certain failure domains.
Considering the unpredictable nature of failure domains in a permissionless system, this behavior is highly undesired since there's no way to prevent those ``bad'' failure domains from emerging.

\sgd{New.}
\lijl{Should clarify that this is a preliminary design.}
\lijl{We can as well give the design principles that guarantees durability, and say this is a possible design instance that follows the principles.}
Guided by the analysis we designed \sys, a permissionless storage system that truly tolerates storage node faults.
\lijl{We need a more precise statement of our property.}
Unlike previous works that all stick on static placement group policies and become vulnerable to corresponding correlated failures, \sys creates placement groups by independently random sampling over all the storage nodes.
Sampling against storage nodes effectively equivalents to simultaneously sampling all failure domains.
Leveraging large erasure coding,
\sgd{need to define ``large'' somewhere before?}
\sys is able to sufficiently sample the failure domains so that the law of large numbers kicks in, ensuring that the times that each failure domain get sampled will be close enough to a known value.
With this we are able to tolerate an individual failure domain by making sure that it is \emph{not} getting sampled for enough times, so that whenever this failure domain fails, all the remaining fragments can still support the recovery.
Because failure domains fail independently to each other, we can also independently tolerate one failure domain by another.
We then determine the erasure coding parameters to tolerate all the failure domains given by a failure generators with high probability.

\sgd{TODO elaboration. (Probably not necessary for a vision paper. Neither the evaluation brief below.)}

\sys is of a practical design.
% We implement \sys on top of libp2p and deploy \sys on 10,000 nodes across 5 continents on AWS.
We implement \sys prototype and deploy \sys on 10,000 nodes across 5 continents on AWS.
\sys shows good performance in both large-scale simulation and physical deployments.
It achieves comparable repair traffic overhead when comparing to a baseline distributed replicated storage system.
Critically, \sys significantly improves fault tolerance to Byzantine participants and targeted adversaries compared to the baseline.
Even when using a coding scheme with small redundancy level, \sys guarantees data durability with more than 33\% Byzantine participants, and more than 10\% of the nodes under targeted attacks.
%\sys performs data object \store and \query with 1.6x-3.1x lower latency compare to existing decentralized storage, while achieving similar scalability over thousands of peers.
\sys performs data object \store with 1.2x-2.2x latency and \query with 0.5-0.9x latency compare to replicated baseline system, while achieving similar scalability over thousands of peers.
\fi
